%==============================================================
%  Figura 5.5 - Estratto del processo iterativo per la progettazione
%  delle parti dei sistemi di comando legate alla sicurezza (SRP/CS)
%  Adattamento in italiano della Figura 5.5 dell'IFA Report 2/2017
%==============================================================
\begin{tikzpicture}[
  font=\normalsize,
  grigio/.style={fill=black!10, align=center, inner sep=4pt},
  numcella/.style={draw=black, fill=white, line width=0.8pt,
                   minimum width=0.7cm, minimum height=0.75cm, inner sep=1pt},
  corpo/.style={draw=black, fill=black!10, line width=0.8pt,
                minimum width=9.0cm, minimum height=0.75cm, inner sep=3pt},
  fdline/.style={draw=black, line width=0.9pt},
  fdar/.style={fdline, -{Triangle[length=6pt,width=6pt]}}
]

%--------------------------------------------------------------
% Ingresso dall'analisi dei rischi
%--------------------------------------------------------------
\node[grigio, minimum width=4.2cm, minimum height=1.3cm] (analisi) at (2.5,-1.05)
  {Dall'analisi dei rischi\\ (EN ISO 12100)};

%--------------------------------------------------------------
% Le tre fasi numerate
%--------------------------------------------------------------
\foreach \n/\y/\testo in {%
  1/-2.65/{Identificazione delle funzioni di sicurezza (SF)},
  2/-3.95/{Specifica delle caratteristiche di ciascuna SF},
  3/-5.25/{Determinazione del PL richiesto (PL\textsubscript{r})}}{
  \node[numcella, anchor=west] (num\n) at (3.3,\y) {\n};
  \node[corpo, anchor=west] (fase\n) at (4.0,\y) {\testo};
}

%--------------------------------------------------------------
% Percorso principale
%--------------------------------------------------------------
\draw[fdar] (analisi.east) -- (5.30,-1.05) -- (5.30,-2.28);
\draw[fdar] (5.30,-3.03) -- (5.30,-3.58);
\draw[fdar] (5.30,-4.33) -- (5.30,-4.88);
\draw[fdar] (5.30,-5.63) -- (5.30,-6.45);

%--------------------------------------------------------------
% "Per ciascuna SF" con graffa
%--------------------------------------------------------------
\node[grigio, minimum width=1.5cm, minimum height=1.6cm] (perogni) at (1.15,-5.25)
  {Per\\ ciascuna\\ SF};
\draw[fdline, decorate, decoration={brace, amplitude=6pt}] (2.60,-5.95) -- (2.60,-4.55);

%--------------------------------------------------------------
% Uscite
%--------------------------------------------------------------
\node[align=center, text width=6.4cm] at (8.15,-7.25)
  {All'implementazione\\ e determinazione del PL\\ (Figura 6.1)};

\draw[fdline] (14.60,-6.45) -- (14.60,-3.95);
\draw[fdar] (14.60,-3.95) -- (13.00,-3.95);
\node[align=center, text width=2.8cm] at (14.60,-7.45)
  {Ritorno\\ se esistono\\ altre SF\\ (Figura 7.1)};

%--------------------------------------------------------------
% Cornice
%--------------------------------------------------------------
\draw[draw=black, line width=1pt]
  ([shift={(-0.45,0.45)}]current bounding box.north west) rectangle
  ([shift={(0.45,-0.45)}]current bounding box.south east);

\end{tikzpicture}
